% !TeX root = main.tex
\documentclass[14pt]{scrartcl}
% !TeX TXS-program:bibliography = txs:///biber
\let\counterwithout\relax
\let\counterwithin\relax
%\usepackage{lmodern}
%\usepackage{minted}
\usepackage{float}
\usepackage{xcolor}

\definecolor{listingrule}{HTML}{B0B0B0}
\definecolor{linenocolor}{HTML}{909090}

\usepackage{extsizes}
\usepackage{subfig}
\usepackage[export]{adjustbox}
\usepackage{tocvsec2}
\usepackage[subfigure]{tocloft}
\usepackage[newfloat,cache=false,outputdir=build]{minted}
\setminted{
  style=refal,
  fontsize=\small,                 % \footnotesize if your column is narrow
  linenos=true,
  numbersep=8pt,                   % gap between numbers and code
  frame=lines,                     % top + bottom rules only
  framesep=2mm,                    % padding inside the frame
  rulecolor=\color{listingrule},
  breaklines=true,
  breakanywhere=false,             % only break at sensible points
  tabsize=2,
  autogobble=true,                 % strip common leading whitespace
  xleftmargin=1em,                 % indent the whole block slightly
  xrightmargin=0pt,
}

\AtBeginEnvironment{figure}{\vspace{0.5cm}}
\AtBeginEnvironment{table}{\vspace{0.5cm}}
\AtBeginEnvironment{listing}{\vspace{0.5cm}}
\AtBeginEnvironment{algorithm}{\vspace{0.5cm}}
\AtBeginEnvironment{minted}{\vspace{-0.5cm}}

% Do not load fancyvrb after minted: minted pulls it in via fvextra, and a
% second load breaks \end{minted} scanning (corrupt \jobname.pyg, FancyVerb errors).
\usepackage{ulem,bm,mathrsfs}
\usepackage{pdflscape}
\usepackage{geometry}
\geometry{a4paper,tmargin=2cm,bmargin=2cm,lmargin=3cm,rmargin=1cm}
\usepackage{amsthm,amsfonts,amsmath,amssymb,amscd}
\usepackage{mathtools}
\IfFileExists{footmisc.sty}{\usepackage[perpage]{footmisc}}{}
\KOMAoptions{fontsize=14pt}

\makeatletter
\def\showfontsize{\f@size{} point}
\makeatother

\makeatletter
\setlength{\@fptop}{0pt}
\makeatother

\usepackage{iftex}
\ifPDFTeX
    \IfFileExists{tempora.sty}{\usepackage{tempora}}{\usepackage{lmodern}}
    \usepackage{cmap}
    \usepackage[T1]{fontenc}
    \usepackage[utf8]{inputenc}
    \usepackage[english,main=russian]{babel}
\else
    \usepackage{fontspec}
    \usepackage{polyglossia}
    \setdefaultlanguage{russian}
    \setotherlanguage{english}
    \defaultfontfeatures{Ligatures=TeX}
    \IfFontExistsTF{Times New Roman}{\setmainfont{Times New Roman}}{\setmainfont{TeX Gyre Termes}}
    \IfFontExistsTF{Times New Roman}{\newfontfamily\cyrillicfont{Times New Roman}}{\newfontfamily\cyrillicfont{TeX Gyre Termes}}
    \IfFontExistsTF{Arial}{\setsansfont{Arial}}{\setsansfont{TeX Gyre Heros}}
    \IfFontExistsTF{Courier New}{\setmonofont{Courier New}}{\setmonofont{DejaVu Sans Mono}}
    \IfFontExistsTF{Arial}{\newfontfamily\cyrillicfontsf{Arial}}{\newfontfamily\cyrillicfontsf{TeX Gyre Heros}}
    \IfFontExistsTF{Courier New}{\newfontfamily\cyrillicfonttt{Courier New}}{\newfontfamily\cyrillicfonttt{DejaVu Sans Mono}}
\fi
\usepackage{indentfirst}
\usepackage{longtable}
\usepackage{array}
\usepackage{booktabs}
\usepackage{soulutf8}
\usepackage{graphicx}
\usepackage{tikz,nicematrix}
\usetikzlibrary{automata, positioning, arrows,patterns, arrows.meta,shadows,shapes,datavisualization,decorations.pathreplacing}
\usetikzlibrary{arrows.meta,patterns.meta,graphs}
\usetikzlibrary{fit, calc}
\usepackage{pgfplots}
\pgfplotsset{compat=1.18}
\usepackage{enumitem}
\usepackage{caption}
\usepackage{setspace}
\AtBeginDocument{\renewcommand{\listingname}{Листинг}}
\usepackage[figure,table,section]{totalcount}
\DeclareTotalCounter{listing}
\IfFileExists{totcount.sty}{\usepackage{totcount}}{}
\IfFileExists{totpages.sty}{\usepackage{totpages}}{}
\linespread{1.3}
\usepackage{hyperref}
\sloppy
\clubpenalty=10000
\widowpenalty=10000

\makeatletter
\let\@@scshape=\scshape
\renewcommand{\scshape}{%
  \ifnum\strcmp{\f@series}{bx}=\z@
    \usefont{T1}{cmr}{bx}{sc}%
  \else
    \ifnum\strcmp{\f@shape}{it}=\z@
      \fontshape{scsl}\selectfont
    \else
      \@@scshape
    \fi
  \fi}
\makeatother

\usepackage{ifthen}
\newcounter{intvl}
\newcounter{otstup}
\newcounter{contnumeq}
\newcounter{contnumfig}
\newcounter{contnumtab}
\newcounter{pgnum}
\newcounter{bibliosel}
\newcounter{chapstyle}
\newcounter{headingdelim}
\newcounter{headingalign}
\newcounter{headingsize}
\newcounter{tabcap}
\newcounter{tablaba}
\newcounter{tabtita}

\setcounter{intvl}{3}
\setcounter{otstup}{0}
\setcounter{contnumeq}{1}
\setcounter{contnumfig}{1}
\setcounter{contnumtab}{1}
\setcounter{pgnum}{0}
\setcounter{bibliosel}{1}
\setcounter{chapstyle}{1}
\setcounter{headingdelim}{1}
\setcounter{headingalign}{0}
\setcounter{headingsize}{0}
\setcounter{tabcap}{0}
\setcounter{tablaba}{2}
\setcounter{tabtita}{1}

\DeclareCaptionLabelSeparator*{emdash}{~--- }
\AddToHook{begindocument/end}{%
  \captionsetup[figure]{labelsep=emdash,font=onehalfspacing,position=bottom}%
  \captionsetup[listing]{position=top,labelsep=emdash}%
}

\ifthenelse{\equal{\thetabcap}{0}}{%
    \newcommand{\tabcapalign}{\raggedright}
}{}
\ifthenelse{\equal{\thetablaba}{0} \AND \equal{\thetabcap}{1}}{%
    \renewcommand{\tabcapalign}{\raggedright}
}{}
\ifthenelse{\equal{\thetablaba}{1} \AND \equal{\thetabcap}{1}}{%
    \renewcommand{\tabcapalign}{\centering}
}{}
\ifthenelse{\equal{\thetablaba}{2} \AND \equal{\thetabcap}{1}}{%
    \renewcommand{\tabcapalign}{\raggedleft}
}{}
\ifthenelse{\equal{\thetabtita}{0} \AND \equal{\thetabcap}{1}}{%
    \newcommand{\tabtitalign}{\raggedright}
}{}
\ifthenelse{\equal{\thetabtita}{1} \AND \equal{\thetabcap}{1}}{%
    \renewcommand{\tabtitalign}{\centering}
}{}
\ifthenelse{\equal{\thetabtita}{2} \AND \equal{\thetabcap}{1}}{%
    \renewcommand{\tabtitalign}{\raggedleft}
}{}

\DeclareCaptionFormat{tablenocaption}{\tabcapalign #1\strut}
\ifthenelse{\equal{\thetabcap}{0}}{%
    \DeclareCaptionFormat{tablecaption}{\tabcapalign #1#2#3}
    \captionsetup[table]{labelsep=emdash}
}{%
    \DeclareCaptionFormat{tablecaption}{\tabcapalign #1#2\par\tabtitalign{#3}}
    \captionsetup[table]{labelsep=space}
}
\captionsetup[table]{format=tablecaption,singlelinecheck=off,font=onehalfspacing,position=top,skip=-5pt}
\DeclareCaptionLabelFormat{continued}{Продолжение таблицы~#2}
\setlength{\belowcaptionskip}{.2cm}
\setlength{\intextsep}{0ex}

\definecolor{linkcolor}{rgb}{0.0,0,0}
\definecolor{citecolor}{rgb}{0,0.0,0}
\definecolor{urlcolor}{rgb}{0,0,0}
\hypersetup{
    linktocpage=true,
    plainpages=true,
    colorlinks,
    linkcolor={linkcolor},
    citecolor={citecolor},
    urlcolor={urlcolor},
    pdflang={ru},
}
\urlstyle{same}
\setlength{\parindent}{2.5em}
\renewcommand{\labelitemi}{\bfseries~{---}}
\setlist{nosep,labelindent=\parindent,leftmargin=*}

\usepackage{ragged2e}
\usepackage[explicit]{titlesec}
\usepackage{xparse}
\pgfdeclarelayer{background}
\pgfdeclarelayer{foreground}
\pgfsetlayers{background,main,foreground}
\makeatletter
\NewDocumentCommand {\getnodedimen} {O{\nodewidth} O{\nodeheight} m} {
  \begin{pgfinterruptboundingbox}
  \begin{scope}[local bounding box=bb@temp]
    \node[inner sep=0pt, fit=(#3)] {};
  \end{scope}
  \path ($(bb@temp.north east)-(bb@temp.south west)$);
  \end{pgfinterruptboundingbox}
  \pgfgetlastxy{#1}{#2}
}
\makeatother
\definecolor{Back_red}{rgb}{1,.45,.34}
\definecolor{Back_blue}{rgb}{0.6,.55,1}
\definecolor{Back_green}{rgb}{0.25,0.9,0.2}
\usepackage{url}
\usepackage{algorithm, algorithmicx}
\usepackage[noend]{algpseudocode}
\usepackage{tabularx}
\usepackage[
    backend=biber,
    bibstyle=gost-numeric,
    sorting=none
]{biblatex}
\addbibresource{biblio.bib}
\titleformat{name=\section,numberless}[block]{\normalfont\Large\centering}{}{0em}{#1}
\titleformat{\section}[block]{\normalfont\Large\bfseries\raggedright}{}{0em}{\thesection\hspace{0.25em}#1}
\titleformat{\subsection}[block]{\normalfont\Large\bfseries\raggedright}{}{0em}{\thesubsection\hspace{0.25em}#1}
\titleformat{\subsubsection}[block]{\normalfont\large\bfseries\raggedright}{}{0em}{\thesubsubsection\hspace{0.25em}#1}
\let\Algorithm\algorithm
\renewcommand\algorithm[1][]{\Algorithm[#1]\setstretch{1.5}}
\usepackage{pifont}
\usepackage{calc}
\makeatletter
\def\p@subsection{}
\def\p@subsubsection{\thesection\,\thesubsection\,}
\makeatother

\usepackage{placeins}
\raggedbottom

\newcommand*\template[1]{\text{<}#1\text{>}}
\newcommand{\enq}[1]{<<#1>>}
\newcommand{\eng}[1]{\begin{english}#1\end{english}}
\newcounter{cTheorem}
\newcounter{cDefinition}
\newcounter{cConsequent}
\newcounter{cExample}
\newcounter{cLemma}
\newcounter{cConjecture}
\newcounter{cProposition}
\newcounter{cRemark}
\theoremstyle{definition}
\newtheorem{Theorem}{Теорема}[cTheorem]
\newtheorem{Definition}{Определение}[cDefinition]
\newtheorem{Consequent}{Следствие}[cConsequent]
\newtheorem{Example}{Пример}[cExample]
\newtheorem{Lemma}{Лемма}[cLemma]
\newtheorem{Conjecture}{Гипотеза}[cConjecture]
\newtheorem{Proposition}{Утверждение}[cProposition]
\newtheorem{Remark}{Замечание}[cRemark]
\newenvironment{Proof}{\begin{proof}}{\end{proof}}
\renewcommand{\theTheorem}{\arabic{Theorem}}
\renewcommand{\theDefinition}{\arabic{Definition}}
\renewcommand{\theConsequent}{\arabic{Consequent}}
\renewcommand{\theExample}{\arabic{Example}}
\renewcommand{\theLemma}{\arabic{Lemma}}
\renewcommand{\theConjecture}{\arabic{Conjecture}}
\renewcommand{\theProposition}{\arabic{Proposition}}
\renewcommand{\theRemark}{\arabic{Remark}}
\NewDocumentCommand{\Newline}{}{\text{\\}}
\newcommand{\sequence}[2]{\ensuremath \left(#1,\ \dots,\ #2\right)}
\definecolor{mygreen}{rgb}{0,0.6,0}
\definecolor{mygray}{rgb}{0.5,0.5,0.5}
\definecolor{mymauve}{rgb}{0.58,0,0.82}
\renewcommand{\listalgorithmname}{Список алгоритмов}
\floatname{algorithm}{Листинг}
\renewcommand{\thealgorithm}{\arabic{algorithm}}
\newcommand{\refAlgo}[1]{(листинг~\ref{#1})}
\newcommand{\refImage}[1]{(рисунок~\ref{#1})}
\newcommand{\inlisting}[1]{в~листинге~\ref{#1}}
\newcommand{\Inlisting}[1]{В~листинге~\ref{#1}}
\newcommand{\onfigure}[1]{на~рисунке~\ref{#1}}
\newcommand{\onfigures}[3]{на~рисунках~\ref{#1}, \ref{#2} и~\ref{#3}}
\newcommand{\Onfigures}[3]{На~рисунках~\ref{#1}, \ref{#2} и~\ref{#3}}
\renewcommand{\theenumi}{\arabic{enumi}.}
\renewcommand{\labelenumi}{\arabic{enumi}.}
\renewcommand{\theenumii}{\arabic{enumii}}
\renewcommand{\labelenumii}{(\arabic{enumii})}
\renewcommand{\theenumiii}{\roman{enumiii}}
\renewcommand{\labelenumiii}{(\roman{enumiii})}
\renewcommand{\labelitemi}{---}
\renewcommand{\labelitemii}{---}
\newcommand{\prog}[1]{{\ttfamily\small#1}}
\newcommand{\anonsection}[1]{\cleardoublepage
\phantomsection
\addcontentsline{toc}{section}{\protect\numberline{}#1}
\section*{#1}\vspace*{2.5ex}
}
\newcommand{\sectionbreak}{\clearpage}
\renewcommand{\cftsecleader}{\cftdotfill{\cftdotsep}}
\renewcommand{\cftsecfont}{\normalfont\large}
\renewcommand{\cftsubsecfont}{\normalfont\large}
\setlength{\cftsecindent}{0pt}
\setlength{\cftsubsecindent}{0pt}
\setlength{\cftsubsubsecindent}{0pt}
\DeclarePairedDelimiter\abs{\lvert}{\rvert}
\DeclarePairedDelimiter\norm{\lVert}{\rVert}
\makeatletter
\DeclareTextCommandDefault{\textvisiblespace}{%
  \mbox{\kern.06em\vrule \@height.3ex}%
  \vbox{\hrule \@width.3em}%
  \hbox{\vrule \@height.3ex}}
\makeatother
\newsavebox{\spacebox}
\begin{lrbox}{\spacebox}
\verb*! !
\end{lrbox}
\newcommand{\aspace}{\usebox{\spacebox}}
\newcommand{\exprLang}{\mathcal{W}_{\text{expr}}}
\newcommand{\patLang}{\mathcal{W}_{\text{pat}}}
\makeatletter
\renewcommand*{\p@subsubsection}{}
\makeatother
\newcommand{\RL}{Рефал5-$\lambda$}
\newcommand{\Lang}{\mathcal{L}}
\DeclareMathOperator{\sent}{\textbf{sent}}
\DeclareMathOperator{\pat}{\textbf{pat}}
\DeclareMathOperator{\abstr}{\textbf{abstr}}
\newcommand{\tuple}[1]{\langle #1 \rangle}
\newcommand{\s}{^\wedge\!}
\newcommand{\e}{\textnormal{\$}}

\def\EAD{\mathcal{E}_{AD}}
\def\AD{AD}

\newcommand{\leqEAD}{\preceq_{\EAD}}
\newcommand{\leEAD}{\prec_{\EAD}}

\newcommand{\litparen}{\textnormal{\texttt{(}}}
\newcommand{\litrep}{\textnormal{\texttt{)}}}
\newcommand{\edgelabel}[1]{%
  \begingroup\footnotesize
  \def\eltxt{#1}%
  \foreach \dx/\dy in {0.55/0,-0.55/0,0/0.55,0/-0.55,0.4/0.4,-0.4/0.4,0.4/-0.4,-0.4/-0.4}{%
    \rlap{\textcolor{white}{\raisebox{\dy pt}{\kern\dx pt\eltxt}}}%
  }%
  \eltxt
  \endgroup
}

\begin{document}
\sloppy

\renewcommand\contentsname{\hfill{\normalfont{СОДЕРЖАНИЕ}}\hfill}
\tableofcontents
\newpage

% 4 стр. титул, кал.план, задание, содержание

\section*{АННОТАЦИЯ} % 1 стр.

\anonsection{ВВЕДЕНИЕ} % 1 стр.

\section{Исследовательская часть}

\subsection{Обзор предметной области} % 2 стр.

Язык программирования Рефал имеет довольно большую историю, начавшуюся 
в 1968 г., когда появилась первая реализация Рефала. Рефал имеет множество диалектов, но все они, по сути, являются надстройками над так называемым базисным Рефалом. Базисный рефал в некотором смысле является самым простым диалектом языка. Программы на базисном Рефале состоят из функций, каждая функция задается парами из шаблона в левой части и выражения в правой.

Единственный тип данных в Рефале --- объектное выражение. Это выражение, 
состоящее из атомарных символов (символ, число) и скобок, при этом скобки правильно вложены. 
Шаблон в левой части функции --- это выражение, в котором, помимо скобок и атомов, могут встречаться переменные. Переменные в шаблоне могут быть трех типов: $e$-переменные, $s$-переменные и $t$-переменные. $e$-переменные могут сопоставляться с любым объектным выражением, $s$-переменные сопоставляются с атомом, $t$-переменные сопоставляются либо с атомом, либо с выражением в скобках. Каждая переменная обладает именем, при этом, если в одном шаблоне есть несколько переменных с одним именем, то при подстановке они заменяются на разные значения.

Основной механизм работы Рефала --- сопоставление с образцом и переписывание выражений. При сопоставлении с образцом, шаблон в левой части функции сопоставляется с объектным выражением в правой части. Если сопоставление прошло успешно, переменным пррисваиваются конкретные значения, затем вычисляется выражение в правой части. 
Сопоставление происходит сверху вниз, если какое-то выражение не сопоставилось ни с одним предложением, это приводит к ошибке времени выполнения.

\Inlisting{lst:IsPal} функция \texttt{IsPal} проверяет, является ли переданное ей выражение, состоящее толлько из атомов, палиндромом.

\begin{listing}[H]
  \caption{Пример функции на базисном Рефале}
  \label{lst:IsPal}
  \begin{minted}{refal}
  IsPal {
    s.x e.x s.x = <IsPal e.x>;
    = True;
    s.oneChar = True;
    e.otherwise = False;
  }
  \end{minted}
\end{listing}

Будучи преимущетсвенно объектом исследования в области суперкомпиляции, Рефал-программы можно исследовать и в контексте статического анализа. 
Статический анализ программ на Рефале представляет интерес по нескольким причинам. Во-первых, специфика языка --- отсутствие явных циклов и присваиваний, опора на сопоставление с образцом и рекурсию --- приводит к тому, что многие классические методы анализа, разработанные для императивных языков, неприменимы напрямую и требуют адаптации. Во-вторых, динамическая природа объектных выражений и наличие переменных, способных сопоставляться с выражениями произвольной длины (в первую очередь $e$-переменных), порождает специфические задачи: например, оценку возможных форм результата функции или определение множества выражений, с которыми может быть вызвана та или иная функция.

\subsection{Общие определения} % 2 стр.

\begin{Definition}[Алфавит]
  Алфавитом $\Sigma$ будем называть конечное множество символов.
\end{Definition}

Будем обозначать алфавиты заглавными греческими буквами.

\begin{Definition}[Связанные со строками понятия]
  Строкой $w = w_1 w_2 \ldots w_n$ будем называть 
  конечную последовательность символов из $\Sigma$.

  Длиной строки $w$ будем называть число $n$, обозначаемое как $|w|$.

  Пустым словом будем называть строку длины 0, обозначаемую как $\varepsilon$.

  $i$-тый символ строки $w$ будем обозначать как $w[i]$.
\end{Definition}

Строки будем обозначать строчными латинскими буквами.

\begin{Definition}[Обращение строки]
Для строки $w = w_1 w_2 \ldots w_n$ будем называть ее обращением строку $w^R = w_n \ldots w_2 w_1$.
Для множества строк $W = \{w_1, w_2, \dots\}$ будем называть его обращением множество $M^R = \{w_n^R, w_2^R, \dots\}$.
\end{Definition}

\begin{Definition}[Итерация Клини]
  Множеством строк $\Sigma^*$ будем называть множество всех конечных строк над алфавитом $\Sigma$.
  $\varepsilon \in \Sigma^*$. Множество $\Sigma^* \setminus \{\varepsilon\}$ обозначим $\Sigma^+$.
\end{Definition}

\begin{Definition}[Сокращенная запись конкатенации]
  Конкатенацией строк $u = u_1 u_2 \ldots u_n$ и $v = v_1 v_2 \ldots v_m$ будем называть 
  строку $u \cdot v = u_1 u_2 \ldots u_n v_1 v_2 \ldots v_m$.

  Если $M \subset \Sigma^*$ --- некоторое множество строк и $w \in \Sigma^*$, 
  то будем обозначать как $Mw$ множество $\{mw \mid  m \in M\}$.
  Аналогично для двух множеств строк $M_1, M_2$ будем обозначать $M_1 M_2$ множество 
  $\{m_1 m_2,\ m_i \in M_i\}$.
\end{Definition}

\begin{Definition}
  Множество $\Sigma^* \setminus W$ будем обозначать $\overline{W}$ и называть дополнением языка $W$.
\end{Definition}

Множества строк будем обозначать заглавными латинскими буквами.

\begin{Definition}[Строковый морфизм]
  Строковый морфизм $\phi\!: \Sigma^* \to \Gamma^*$ --- это функция, которая каждому слову $w \in \Sigma^*$ 
  ставит в соответствие слово $\phi(w) \in \Gamma^*$, сохраняющая конкатенацию, т.е. $\forall w_1, w_2 \in \Sigma^* \ \phi(w_1 w_2) = \phi(w_1) \phi(w_2)$. Строковый морфизм однозначно задается набором значений на символах алфавита $\Sigma$. Если для некоторого $\sigma \in \Sigma$ не задано явным образом отображение, будем считать, что $\sigma$ отображается в себя. В таком случае будем опускать этот факт в записи областей определения и значений морфизма.
\end{Definition}

Все морфизмы в данной работе --- строковые, поэтому будем опускать слово "строковый".
Будем обозначать морфизмы строчными греческими буквами.

\begin{Example}
  Пусть $\Sigma = \Gamma = \{a, b, c\}$. Морфизм, меняющий местами $a$ и $b$, можно задать следующим образом:
  $\phi\!: \{a, b\} \to \{a, b\}$, при этом 
  \begin{align*}
    & \phi(a) \mapsto b; \\
    & \phi(b) \mapsto a.
  \end{align*}
  В таком способе записи неявно задано $\phi(c) = c$.
\end{Example}

\begin{Definition}[синтаксическая конгруэнция]
Синтаксической конгруэнцией языка $L$ будем называть отношение эквивалетности $\sim_L$ такое, что
$$
u \sim_L v \iff \forall x, y \in \Sigma^* \ (xuy \in L \iff xvy \in L)
$$
\end{Definition}

Класс эквивалентности некоторого $x$ в фактор-структуре будем обозначать, как $[x]$

\begin{Definition}[Синтаксический моноид]
Фактор-множество $\Sigma^* /\! \sim_L$ с операцией конкатенации ($[x][y] = [xy]$) --- 
синтаксический моноид языка $L$.
\end{Definition}

\begin{Definition}[синтаксическая полугруппа]
  Фактор-множество $\Sigma^+ /\! \sim_L$ с операцией конкатенации ($[x][y] = [xy]$) --- 
  синтаксическая полугруппа языка $L$.
\end{Definition}

Для $n \in \mathbb{N}$ примем $\overline{1,n} = \{1, 2, \ldots, n\}$.

\subsection{Перекрытия} % 2 стр.

\begin{Definition}[Язык правильно вложенных слов]
  \label{Def::PLIS}
  Пусть $\Sigma_{\text{call}}$, $\Sigma_{\text{ret}}$ и $\Sigma_{\text{int}}$ --- непересекающиеся алфавиты. 
  Язык правильно вложенных слов $\mathcal{W}$ --- это наименьшее множество, 
  которое генерируется следующими конструкторами:
  \begin{align*}
  \varepsilon &\in \mathcal{W} \quad \text{(пустое слово)} \\
  \gamma &\in \mathcal{W} \quad \text{для каждого } \gamma \in \Sigma_{\text{int}} \quad \text{(внутренний символ)} \\
  w_1 \cdot w_2 &\in \mathcal{W} \quad \text{если } w_1, w_2 \in \mathcal{W}  \\
  \langle_c \; w \; \rangle_r &\in \mathcal{W} \quad \text{если } w \in \mathcal{W},\; \langle_c \in \Sigma_{\text{call}},\; \rangle_r \in \Sigma_{\text{ret}} 
  \end{align*}
\end{Definition}

Выражением будем называть слова языка правильно вложенных слов $\exprLang$, 
порожденного алфавитами $\Sigma_{\text{call}} = \{\litparen\}$, 
$\Sigma_{\text{ret}} = \{\litrep\}$ и $\Sigma_{\text{int}} = \Delta$, 
где $\Delta$ --- счетное множество символов.
Абстрактным шаблоном без повторных переменных будем называть слова языка правильно вложенных слов $\patLang$, 
порожденного алфавитами $\Sigma_{\text{call}} = \{\litparen\}$, 
$\Sigma_{\text{ret}} = \{\litrep\}$ и $\Sigma_{\text{int}} = \{e, s, t\}$.

\begin{Remark}
Формальное описание абстрактного шаблона без повторных переменных соответсвует таким шаблонам языка Рефал, в которых все переменные имеют различные имена и атомарные символы не встречаются вовсе.
Например, для диалекта \RL, поддерживающего безымянные переменные,
из абстрактного шаблона можно получить обычный, заменив все символы $e$ ($s, t$)
абстрактного шаблона на символы \texttt{e.\_} (\texttt{s.\_},\texttt{t.\_}) шаблона \RL.
\end{Remark}

Введем также понятие шаблона с пронумерованными переменными: каждому шаблону $p \in \patLang$ 
поставим в соответствие слово над алфавитом $\Delta_\mathbb{N} = \{\litparen, \litrep\} \cup \{e, s, t\} \times \mathbb{N}$ такое,
что все вхождения символов $e, s, t$ в $p$ в шаблоне с пронумерованными переменными 
заменяются на буквы из $\{e, s, t\} \times \mathbb{N}$, причем $i$-тому вхождению 
символа $e$ в $p$ соответствует буква $(e, i)$ (обозначим $e_i$). Аналогично для $s$ и $t$. 
Скобки $\litparen, \litrep$ остаются на месте. Для шаблона $p$ будем обозначать соответствующий ему 
шаблон с пронумерованными переменными $p_\mathbb{N}$

Этот формализм нужен лишь затем, чтобы формализовать понятие сопоставления с образцом.

Соответствие переменных множествам их допустимых значений было приведено 
в обзоре базисного Рефала, теперь же формализуем это соответствие с учетом, что рассматриваются только абстрактные шаблоны без повторных переменных.
Будем говорить, что выражение $w$ \textit{допускает сопоставление} с образцом $p$, 
если существует строковый морфизм $\phi\!: \Delta_\mathbb{N}^* \to \exprLang$, действующий следующим образом:
\begin{align*}
  & \phi((s, i)) \mapsto \gamma \in \Sigma_{\text{int}} \ \forall i \in \mathbb{N}; \\
  & \phi((t, i)) \mapsto \gamma \in \Sigma_{\text{int}} \text{ или }
    \phi((t, i)) \mapsto \litparen\, v \,\litrep,\ v \in \exprLang \ \forall i \in \mathbb{N}, \\
\end{align*}
при этом $\phi(p_\mathbb{N}) = w$.

\begin{Definition}[Язык шаблона]
Языком шаблона $p \in \patLang$ будем называть множество всех 
выражений $w \in \exprLang$, которые допускают сопоставление с $p$.
Обозначать это множество будем как $\Lang(p)$.
\end{Definition}

Каждая функция Рефала задается последовательностью предложений, каждое из которых
в свою очередь содержит шаблон в левой части.
Для некоторой функции $f$, содержащей $n$ предложений, будем обозначать 
$i$-тый по счету (считая сверху вниз) шаблон предложения
как $f.p_i$. Последовательность всех шаблонов функции 
$\tuple{f.p_1, f.p_2, \ldots, f.p_n}$ будем обозначать как $\pat f$.

\begin{Definition}[Перекрытие]
Будем говорить, что шаблон $p$ \textit{перекрывается} или 
\textit{экранируется} шаблонами $p_1, p_2, \ldots, p_k$, если 
$$
\Lang(p) \subseteq \bigcup_{i=1}^k \Lang(p_i).
$$
Будем говорить, что перекрытие (экранирование) \textit{множественное}, если $k > 1$.
\end{Definition}

Итак, после введения всех необходимых понятий, 
мы можем перейти к определению задачи распознавания перекрытия:
\begin{Definition}[Задача распознавания перекрытия]
Задача распознавания перекрытия в функции $f$ с $n$ предложениями 
заключается в том, чтобы для каждого шаблона $p_i \in \pat f,\ i \in \overline{1,n}$ 
найти минимальное по включению множество 
$I = \{i_1, i_2, \ldots, i_k\} \subseteq \overline{1,i-1}$ такое, что 
$p_i$ перекрывается шаблонами $f.p_{i_1}, f.p_{i_2}, \ldots, f.p_{i_k}$, 
или показать, что такого множества не существует.
\end{Definition}

Отметим, что решение задачи распознавания перекрытия неоднозначно, так 
как для одного и того же шаблона можно найти несколько минимальных по 
включению множеств индексов таких, что соответствующие им шаблоны 
перекрывают искомый. Пример такого случая приведён \inlisting{lst:ambiguous-overlapping}.
Предложение~5 перекрывается предложениями~3 и~4, а также предложением~1. 
При этом множества $\{3, 4\}$ и $\{1\}$ несравнимы по включению.

\begin{listing}[H]
\caption{Пример неоднозначного экранирования предложений Рефал-функции}
\label{lst:ambiguous-overlapping}
\begin{minted}{refal}
F {
    (t._ t._) e._ = ...;          /* 1 */
    e._ (s._ e._ s._) e._ = ...;  /* 2 */
    e._ s._ (t._) = ...;          /* 3 */
    e._ (e._) (t._) = ...;        /* 4 */
    (t._ t._) t._ (t._) = ...;    /* 5 */
}
\end{minted}
\end{listing}

Забегая вперед, отметим, что в текущей постановке задача распознавания перекрытия 
алгоритмически труднее, чем ее упрощенный вариант:
\begin{Definition}[Упрощенная задача распознавания перекрытия]
Упрощенная задача распознавания перекрытия в функции $f$
заключается в том, чтобы для каждого шаблона $p_i \in \pat f$ определить, 
экранируется ли он шаблонами $p_1, p_2, \ldots, p_{i-1}$.
\end{Definition}

Искомое определению задачи распознавания перекрытия будем называть теперь полным или точным. 

\subsection{Регулярное представления шаблонов} % 0.5 стр. 

Регулярное представление шаблона представляет особый интерес, поскольку, в сравнении 
с контекстно-свободными языками (в частности, языками с магазинных автоматов, управляемых входом), 
коими являются $\exprLang$ и $\patLang$, регулярные языки в каком-то смысле более просты для анализа. 

\subsubsection{Каскадное произведение автоматов} % 3.5 стр.

Рассмотрим основную теоретическую конструкцию, 
позволяющую извлекать регулярное представление шаблонов --- каскадное произведение автоматов:

\begin{Definition}[Каскадное произведение автоматов]
  Пусть $A_1 = \langle Q_1, \Sigma, \delta_1, q_{0_1}, F_1 \rangle$ и
  $A_2 = \langle Q_2, \Sigma \times Q_1, \delta_2, q_{0_2}, F_2 \rangle$ --- детерминированные конечные автоматы.
  \emph{Каскадным произведением} $A_1 \circ A_2$ называется ДКА $\langle Q, \Sigma, \delta, q_0, F \rangle$, где
  \begin{align*}
    Q &= Q_1 \times Q_2, \\
    q_0 &= (q_{0_1}, q_{0_2}), \\
    F &= F_1 \times F_2, \\
    \delta((q_1, q_2), a) &= \bigl( \delta_1(q_1, a),\; \delta_2(q_2, (a, q_1)) \bigr).
  \end{align*}
\end{Definition}

Иногда будем просто говорить о каскаде (двух и более автоматов), подразумевая каскадное произведение автоматов.

Определение, данное выше, является вариантом прямого произведения автоматов. Как и многие другие операции над ДКА, использующие прямое произведение (например, пересечение и объединение), каскадное произведение можно описать как систему из двух автоматов, 
работающих параллельно на одном и том же входном слове. В случае каскада двух автоматов между ними проявляется односторонняя зависимость:
оба автомата синхронно читают посимвольно входное слово $w \in \Sigma^*$, при этом первый автомат 
$A$ работает автономно над алфавитом $\Sigma$, а второй автомат, помимо символа $w[i]$, 
использует текущее (до перехода по $w[i]$) состояние первого автомата $q$, т.е. работает над алфавитом $\Sigma \times Q_1$.

Основная идея заключается в делегировании распознования вложенных подшаблоннов элементам каскада.

Введем искомый алфавит $\Gamma_0 = \{b, s, \litparen, \litrep\}$. $b$ является мета-символом, 
соответствующим словам $\exprLang \ni w = \litparen\, v \,\litrep,\ v \in \exprLang$, $s$ также является мета-символом, 
соответствующим словам $\exprLang \ni w = \gamma \in \Sigma_{\text{int}}$.
Будем рассматривать регулярные выражения $r$ над алфавитом $\Gamma_0$ и ДКА $A$ над этим же алфавитом. 
Будем обозночать множество слов $\{w \mid w \in \Gamma_0^*\}$, принадлежащие языку $r$ ($A$), как $R(r)$ ($R(r)$).
Отклонение от стандартного обозначения через $\Lang$ сделано во избежании путаницы: каждое слово языка $r$ ($A$) 
в силу смысла символов $\gamma \in \Gamma_0$ само задает некоторый язык слов-выражений из $\exprLang$.
Для слова $u \in R(r)$ поставим ему в соответствие слово $u_\mathbb{N}$ над 
$\Gamma_\mathbb{N} = \{\litparen, \litrep\} \cup \{b, s\} \times \mathbb{N}$ с пронумерованными переменными, 
аналогично тому, как это было сделано для шаблонов $p \in \patLang$ и $p_\mathbb{N}$.
Будем говорить, что $w$ допускает сопоставление с $r$, если
$$
\exists u \in R(r)\ \exists \phi\!: \{b, s\} \times \mathbb{N} \to \exprLang \mid \phi(u_\mathbb{N}) = w,
$$
при этом
\begin{align*}
& \phi((b, i)) \mapsto \litparen\, w \,\litrep \in \exprLang \forall i \in \mathbb{N}; \\
& \phi((s, i)) \mapsto \gamma \in \Sigma_{\text{int}} \forall i \in \mathbb{N};
\end{align*}

Множество слов $w \in \exprLang$, допускающих сопоставление с $r$, будем обозначать как $\Lang(r)$.
Любой шаблон $p \in \patLang$ можно представить в виде регулярного выражения $r$ 
над алфавитом $\Gamma$ так, что $\Lang(p) = \Lang(r)$.
Данный факт следует непосредственно из определения самих $e, t, s$-символов, 
ниже представлено соответствие между символами и регулярными выражениями:
\begin{itemize}
\item $e$-символ соответствует выражению $(b|s)^*$;
\item $t$-символ соответствует выражению $(b|s)$;
\item $s$-символ соответствует выражению $s$.
\end{itemize}

\subsubsection{Переход от каскада к ДКА} % 2.5 стр.

\subsection{Распознавания перекрытий с помощью ДКА шаблонов} % 1 стр.

\subsection{Иные применения ДКА шаблонов} % 2 стр.
% пара примеров (простой и не очень), где выводится уточнение 

\section{Конструкторская часть}

\subsection{Общая архитектура программного комплекса} % 2 стр.
% диаграммка: справа лексер -> синтаксический анализатор -> семантический анализатор,
% слева обертка CLI - API

\subsubsection{Обработка ошибок} % 1 стр.

\subsection{Архитектура подсистемы семантического анализатора} % 1 стр.

\subsubsection{Вспомогательные подмодули семантического анализатора} % 2-3 стр.

\subsubsection{Алгоритм приведения шаблонов функции к ДКА} % 3 стр.
% в т.ч. замечание о производительности в случае детального анализа

\subsection{Алгоритмы извлечения документации к шаблонам функции} % 3 стр.

\subsection{Пользовательский интерфейс} % 2 стр.

\section{Технологическая часть}

\subsection{Обзор используемых технологий} % 2 стр.

\subsection{Реализация модуля для работы с абстрактными паттернами} % 2 стр.

\subsection{Реализация модуля для работы с автоматами} % 2 стр.

\subsection{Реализация алгоритма обнаружения перекрытий} % 2 стр.

\subsection{Реализация алгоритма извлечения документации} % 3 стр.

\subsection{Руководство пользователя} % 2 стр.
% что и как установить
% git
% что как запускат

\section{Аналитическая часть}
% WiP

\anonsection{ЗАКЛЮЧЕНИЕ} % 1 стр.

\anonsection{СПИСОК ЛИТЕРАТУРЫ} % 1 стр.

\printbibliography[title={Список литературы},heading=bibintoc]

\end{document}
